% Part: sets-functions-relations
% Chapter: sets
% Section: pairs-and-products

\documentclass[../../../include/open-logic-section]{subfiles}

\begin{document}

\olfileid[id]{sfr}{set}{pai}
\olsection{Pasangan, Tupel, dan Hasil Kali Kartesius}

\begin{explain}
Dari ekstensionalitas diperoleh bahwa anggota sebuah himpunan
tidak mempunyai urutan. Karena itu, jika kita ingin merepresentasikan
urutan, kita menggunakan \emph{pasangan terurut} $\tuple{x, y}$.
Dalam pasangan tak terurut $\{x, y\}$, urutan tidak berpengaruh:
$\{x, y\} = \{y, x\}$. Dalam pasangan terurut, urutan berpengaruh:
jika $x \neq y$, maka $\tuple{x, y} \neq \tuple{y, x}$.

Bagaimana kita seharusnya memahami pasangan terurut dalam teori
himpunan? Yang sangat penting, kita ingin mempertahankan gagasan bahwa
dua pasangan terurut identik jika dan hanya jika keduanya mempunyai
anggota pertama yang sama dan anggota kedua yang sama, yaitu:
\[
  \tuple{a, b}= \tuple{c, d}\text{ jika dan hanya jika }a = c
  \text{ dan }b=d.
\]
Kita dapat mendefinisikan pasangan terurut dalam teori himpunan dengan
definisi Wiener--Kuratowski.
\end{explain}

\begin{defn}[Pasangan terurut]\ollabel{wienerkuratowski}
	$\tuple{a, b} = \{\{a\}, \{a, b\}\}$.
\end{defn}

\begin{prob}
	Dengan menggunakan \olref[sfr][set][pai]{wienerkuratowski}, buktikan
	bahwa $\tuple{a, b}= \tuple{c, d}$ jika dan hanya jika $a = c$ dan
	$b=d$.
\end{prob}

\begin{explain}
Setelah menetapkan definisi pasangan terurut, kita dapat menggunakannya
untuk mendefinisikan himpunan lain. Sebagai contoh, kadang-kadang kita
juga memerlukan barisan terurut yang terdiri atas lebih dari dua objek,
misalnya \emph{tripel} $\tuple{x, y, z}$, \emph{kuadruplet}
$\tuple{x, y, z, u}$, dan seterusnya. Kita dapat memandang tripel
sebagai pasangan terurut khusus yang anggota pertamanya sendiri
merupakan pasangan terurut: $\tuple{x, y, z}$ adalah
$\tuple{\tuple{x, y},z}$. Hal yang sama berlaku bagi kuadruplet:
$\tuple{x,y,z,u}$ adalah $\tuple{\tuple{\tuple{x,y},z},u}$, dan
seterusnya. Secara umum, kita berbicara tentang \emph{tupel-$n$
terurut} $\tuple{x_1, \dots, x_n}$.

Himpunan tertentu yang terdiri atas pasangan terurut atau tupel-$n$
terurut lainnya akan berguna.
\end{explain}

\begin{defn}[Hasil kali Kartesius]
Untuk himpunan $A$ dan $B$, \emph{hasil kali Kartesius} $A \times B$
didefinisikan oleh
\[
  A \times B = \Setabs{\tuple{x, y}}{x \in A \text{ dan } y \in B}.
\]
\end{defn}

\begin{ex}
Jika $A = \{0, 1\}$ dan $B = \{1, a, b\}$, maka hasil kali keduanya
adalah
\[
A \times B = \{ \tuple{0, 1}, \tuple{0, a}, \tuple{0, b},
    \tuple{1, 1}, \tuple{1, a}, \tuple{1, b} \}.
\]
\end{ex}

\begin{ex}
Jika $A$ adalah himpunan, hasil kali $A$ dengan dirinya sendiri,
$A \times A$, juga ditulis~$A^2$. Ini adalah himpunan \emph{semua}
pasangan $\tuple{x, y}$ dengan $x, y \in A$. Himpunan semua tripel
$\tuple{x, y, z}$ adalah $A^3$, dan seterusnya. Kita dapat memberikan
definisi rekursif:
\begin{align*}
  A^1 & = A\\
  A^{k+1} & = A^k \times A
\end{align*}
\end{ex}

\begin{prob}
Tuliskan semua !!{element}s dari $\{1, 2, 3\}^3$.
\end{prob}

\begin{prop}\ollabel{cardnmprod}
Jika $A$ mempunyai $n$ !!{element}s dan $B$ mempunyai $m$
!!{element}s, maka $A \times B$ mempunyai $n\cdot m$ anggota.
\end{prop}

\begin{proof}
Untuk setiap !!{element}~$x$ dalam~$A$, terdapat $m$ !!{element}s
berbentuk $\tuple{x, y} \in A \times B$. Misalkan
$B_x = \Setabs{\tuple{x, y}}{y \in B}$. Karena jika $x_1 \neq x_2$,
maka $\tuple{x_1, y} \neq \tuple{x_2, y}$, diperoleh
$B_{x_1} \cap B_{x_2} = \emptyset$. Akan tetapi, jika
$A = \{x_1, \dots, x_n\}$, maka
$A \times B = B_{x_1} \cup \dots \cup B_{x_n}$, sehingga himpunan
ini mempunyai $n\cdot m$ !!{element}s.

Untuk memvisualisasikannya, susun !!{element}s dari~$A \times B$
dalam sebuah kisi:
\[
\begin{array}{rcccc}
  B_{x_1} = & \{\tuple{x_1, y_1} & \tuple{x_1, y_2} & \dots & \tuple{x_1, y_m}\}\\
  B_{x_2} = & \{\tuple{x_2, y_1} & \tuple{x_2, y_2} & \dots & \tuple{x_2, y_m}\}\\
  \vdots & & \vdots\\
  B_{x_n} = & \{\tuple{x_n, y_1} & \tuple{x_n, y_2} & \dots & \tuple{x_n, y_m}\}
\end{array}
\]
Karena semua $x_i$ berbeda dan semua $y_j$ berbeda, tidak ada dua
pasangan dalam kisi ini yang sama, dan jumlah pasangan itu adalah
$n\cdot m$.
\end{proof}

\begin{prob}
Tunjukkan, dengan induksi pada~$k$, bahwa untuk setiap $k \ge 1$, jika
$A$ mempunyai $n$ !!{element}s, maka $A^k$ mempunyai $n^k$
!!{element}s.
\end{prob}

\begin{ex}
Jika $A$ adalah himpunan, sebuah \emph{kata} atas~$A$ adalah sebarang
barisan !!{element}s dari~$A$. Sebuah barisan dapat dipandang sebagai
tupel-$n$ dari !!{element}s dalam~$A$. Sebagai contoh, jika
$A = \{a, b, c\}$, maka barisan ``$bac$'' dapat dipandang sebagai
tripel~$\tuple{b, a, c}$. Kata, yaitu barisan simbol, sangat penting
dalam ilmu komputer. Berdasarkan konvensi, kita menghitung
!!{element}s dari~$A$ sebagai barisan dengan panjang~$1$, dan
$\emptyset$ sebagai barisan dengan panjang~$0$. Himpunan \emph{semua}
kata atas~$A$ kemudian adalah
\[
A^* = \{\emptyset\} \cup A \cup A^2 \cup A^3 \cup \dots
\]
\end{ex}

\end{document}
